%% kapa: exam=2 round=29 year=2018 period=1 subject=appraisal_practice subject_ko=감정평가실무 seq=3 number=3 points=20
%% source: Q-Net 감정평가사 기출문제 (cst00309) artlSeq=5203398 file=raw/1008/5203398/1교시 감정평가실무.pdf
%% conversion: kordoc 4.24.51 -> tools/gen2.py (auto); review=proofread-sonnet; flags=none
\begin{longquestion}{【문제 3】}
베트남 북동해역에서 석유시추용으로 운용되던 플랜트(선박)인 ‘스타호’는 경제성 저하 및 노후화로 ‘비운영 폐선’으로의 매각을 추진 중이며, 현재 싱가포르 외항에 정박 중이다. 소유자인 코리아석유공사는 2018. 06. 30. 기준의 유리한 매각방식을 결정하기 위한 자문을 감정평가사 甲에게 구하였다. 관련 법규 및 감정평가이론을 참작하고 제시된 자료를 활용하여 다음의 물음에 답하시오. (20점) \par
\medskip\noindent 물음 1) 해체처분가격의 성격을 약술하고, 전체를 해체처분가격으로 평가할 경우, 산출 가능한 시산가액을 매각처 별로 산정하시오. (10점) \par
\medskip\noindent 물음 2) 재사용이 가능한 기관 및 저장품은 분리하여 매각할 경우의 전체 시산가액을 산정한 후, 물음 1)과 비교하여 가장 유리한 매각방식을 결정하시오. (10점) \par
< 자료 1 > ‘스타호’의 개요 \par
1. 종류: 부선 \par
2. 선질: 강 \par
3. 조선자: 울산조선(주) \par
4. 진수일: 1990. 06. 30. \par
5. 길이: 75미터 \par
6. 너비: 60미터 \par
7. 깊이: 8미터 \par
8. 총톤수: 10,000톤 \par
9. 재화중량(dead weight): 13,000톤 \par
10. 경하중량(light weight): 15,000톤 \par
11. 기관: 디젤엔진(2,000hp) 2대 탑재되어 있고, 중량은 총 100톤으로 조사됨 \par
12. 저장품: 선박에 탑재된 수리용 신품의 부속장비로 중량은 총 900톤으로 조사됨 \par
< 자료 2 > 가격조사 사항 \par
1. 통상 선박의 해체처분은 정상운영 장비가 포함된 경하중량을 기준으로 거래되는 관행이며, 대형 선박 또는 플랜트의 해체 조선소는 파키스탄 및 한국에 소재함 \par
2. 기준시점 현재 현지 인도조건의 scrap(고철) 매입단가는 파키스탄의 경우 톤당 260,000원 수준이고, 한국의 경우 톤당 240,000원 수준인 것으로 조사됨 \par
3. 한편, 싱가포르 소재 대형선박 및 플랜트 관련 에이전트는 톤당 200,000원 수준에서 즉시 매입의사를 밝히고 있음 \par
4. 본건은 자력항행이 불가능한 부선으로 현지 인도조건에 따른 운송비(보험료 포함)는 파키스탄의 경우 9억원, 한국의 경우 6억원이 소요되는 것으로 조사되고, 싱가포르 현지매각의 경우 매수자가 모든 부대비용을 부담하는 조건임 \par
5. 재사용 가능 부분의 분리매각의 경우, 원매자 탐색 및 분리작업 기간에 4개월이 소요되고, 이에 따른 매월 정박료 및 대기비용으로 월간 2억원의 부담이 예상되나, 분리에 따른 작업 직접비용은 매수자 부담이며, 잔여 scrap(고철)은 싱가포르 현지에서 매각 예정임 \par
< 자료 3 > 재조달원가 \par
1. 기준시점에서 기관의 재조달원가는 마력(hp)당 300,000원인 것으로 조사됨 \par
2. 저장품은 미사용품으로 취득가격은 50억원이며, 이를 재조달원가로 할 수 있음 \par
3. 선체 및 의장품은 노후화로 본래용도로의 재사용은 불가능할 것으로 판단함 \par
< 자료 4 > 내용연수 및 잔존가치율 등 \par
1. 기관의 내용연수는 20년이고, 잔존가치율은 10 \%이나, 매년 정기적 유지보수로 경제적 측면의 잔존내용연수가 5년 정도 남아있는 것으로 조사되며, 정률법에 의한 감가수정을 함 \par
2. 저장품의 내용연수는 10년이고, 잔존가치율은 20 \%이며, 미사용 신품이지만 원매자가 제한되어 있어 잔존가치 정도에서 거래가 가능함 \par
< 자료 5 > 정률법에 의한 잔존가치율 표 \par
{\small\centering
\begin{tabular}{|c|c|c|}
\hline
\makecell{구분} & \multicolumn{2}{|c|}{\makecell{잔존가치율(10 \%)}} \\
\hline
\makecell{내용연수\newline{}경과연수} & \makecell{15년} & \makecell{20년} \\
\hline
\makecell{1} & \makecell{14/0.858} & \makecell{19/0.891} \\
\hline
\makecell{2} & \makecell{13/0.736} & \makecell{18/0.794} \\
\hline
\makecell{3} & \makecell{12/0.631} & \makecell{17/0.708} \\
\hline
\makecell{4} & \makecell{11/0.541} & \makecell{16/0.631} \\
\hline
\makecell{5} & \makecell{10/0.464} & \makecell{15/0.562} \\
\hline
\makecell{6} & \makecell{9/0.398} & \makecell{14/0.501} \\
\hline
\makecell{7} & \makecell{8/0.341} & \makecell{13/0.447} \\
\hline
\makecell{8} & \makecell{7/0.293} & \makecell{12/0.398} \\
\hline
\makecell{9} & \makecell{6/0.251} & \makecell{11/0.355} \\
\hline
\makecell{10} & \makecell{5/0.215} & \makecell{10/0.316} \\
\hline
\makecell{11} & \makecell{4/0.185} & \makecell{9/0.282} \\
\hline
\makecell{12} & \makecell{3/0.158} & \makecell{8/0.251} \\
\hline
\makecell{13} & \makecell{2/0.136} & \makecell{7/0.224} \\
\hline
\makecell{14} & \makecell{1/0.117} & \makecell{6/0.200} \\
\hline
\makecell{15} & \makecell{0.100} & \makecell{5/0.178} \\
\hline
\makecell{16} &  & \makecell{4/0.158} \\
\hline
\makecell{17} &  & \makecell{3/0.141} \\
\hline
\makecell{18} &  & \makecell{2/0.126} \\
\hline
\makecell{19} &  & \makecell{1/0.112} \\
\hline
\makecell{20} &  & \makecell{0.100} \\
\hline
\end{tabular}\par}

\end{longquestion}
